\section{Phonon-Phonon SSE with decomposed FC3}
\label{sec:phph_decomp}
% ----------------------------------------------------------------------

The sparsity layers of \cref{ssec:phph_sparsity} reduce the per-pair
ring contraction from $\cO(\NBS^{4})$ to $\cO(\NBS\,C_a^{3})$. The
factor $C_a^{3}$ remains the dominant cost in
\cref{eq:cost_phph_total}. We now show that any decomposition of the FC3
in \cref{sub:fc3_compression} reduces this further to
$\cO(R\,\NBS\,r_c\,C_a + R^{2}\NBS)$ by collapsing the four-cycle of
the ring contraction into a single Gram matrix on the propagator.
The reduction is independent of which decomposition is used in the
sense that all CP-family decompositions (CP, INDSCAL, Waring, PCP) produce
the same kernel structure. The symmetry constraints distinguishing
them affect only which factor vectors enter the Gram product and how
much it can be symmetrised.

We derive this under two assumptions that hold for the SSE of a single
SCBA iteration on a single FC3 dataset. First, both vertices of the
sunset diagram (\cref{fig:sunset_diagram}) are evaluated against the
same $\Phi^{(3)}$: $\Phi_L \equiv \Phi_R \equiv \Phi$. Second, the
two internal propagator lines are restrictions of the same Green's
function $g(\tau)$ to (possibly different) block-row supports.

The section is organised as follows.
\Cref{ssec:phph_decomp_setup} fixes the notation and the abstract
separable form. \Cref{ssec:phph_decomp_kernel} derives the factorised
kernel and identifies the algebraic simplification that distinguishes
the symmetric ans\"atze from plain CP.
\Cref{ssec:phph_decomp_ansatze} specialises to each ansatz of
\cref{sub:fc3_compression}. \Cref{ssec:phph_decomp_blocks} traces
through the three block-sparsity layers of \cref{ssec:phph_sparsity}
under decomposition. \Cref{ssec:phph_decomp_compute} gives a concrete
three-step algorithm and its cost. \Cref{ssec:phph_decomp_qspace}
addresses the partial Fourier transform.
\Cref{ssec:phph_decomp_dist} reworks the distributed pipeline of
\cref{ssec:phph_distribution}.

% ----------------------------------------------------------------------
\subsection{Setup}
\label{ssec:phph_decomp_setup}
% ----------------------------------------------------------------------

Every decomposition considered in \cref{sub:fc3_compression} writes the FC3 vertex as a sum of
separable terms. The most general form is
\begin{equation}\label{eq:fc3_separable_general}
  \Phi_{x y_1 y_2}
  \;\approx\;
  \sum_{r=1}^{R} \lambda_r\;\alpha_r(x)\;\beta_r(y_1)\;\gamma_r(y_2),
\end{equation}
with weights $\lambda_r\in\mathbb{R}$ and factor vectors
$\alpha_r,\beta_r,\gamma_r\in\mathbb{R}^{\ndof}$. The symmetry constraints are
\begin{equation}\label{eq:symmetry_constraints}
\begin{aligned}
  \text{CP:}      \quad& \text{no constraint,}\\
  \text{INDSCAL:} \quad& \beta_r=\gamma_r,\\
  \text{Waring:}  \quad& \alpha_r=\beta_r=\gamma_r,\\
  \text{Tucker:}  \quad& \text{$\beta,\gamma$ taken from a common $R_2$-dim basis,}\\
                   & \text{$\alpha$ from a separate $R_1$-dim basis, plus a small core.}
\end{aligned}
\end{equation}
These constraints reduce parameter count
(\cref{eq:fc3_matsvd_params,eq:fc3_tucker_params,eq:fc3_cp_params,eq:fc3_indscal_params,eq:fc3_waring_params,eq:fc3_pcp_params})
and enforce permutation symmetry of
the reconstructed tensor in $(y_1,y_2)$ (INDSCAL, Waring) or all
three legs (Waring, PCP). PCP is treated separately in
\cref{ssec:phph_decomp_ansatze} because its $6$-term explicit
symmetrisation interacts non-trivially with the kernel.

We collect the factor vectors as columns of three matrices
$A,B,C \in \mathbb{R}^{\ndof\times R}$ with
$A_{x,r}\!=\!\alpha_r(x)$, $B_{y,r}\!=\!\beta_r(y)$,
$C_{y,r}\!=\!\gamma_r(y)$, and the weights as a diagonal matrix
$\Lambda = \diag(\lambda_1,\dots,\lambda_R)$. Under the constraints
\cref{eq:symmetry_constraints}, $B=C$ for INDSCAL and $A=B=C\equiv U$
for Waring. The propagator
% REVIEW: time-domain greater/lesser GF is complex, not real
$g\equiv g(\tau)\in\mathbb{C}^{\ndof\times\ndof}$
is the time-domain Green's function obtained from Stage~1 of
\cref{ssec:phph_distribution}. We will not yet distinguish lesser,
greater, or retarded components, the derivation applies unchanged to
all three.

% ----------------------------------------------------------------------
\subsection{Factorised kernel}
\label{ssec:phph_decomp_kernel}
% ----------------------------------------------------------------------

The per-pair ring \cref{eq:inner_contraction} is
\begin{equation}
  \mathcal{B}_{xx'}
  = \sum_{y_1 y_2 y_3 y_4}
    \Phi(x,y_1,y_2)\,g(y_1,y_4)\,g(y_2,y_3)\,
    \Phi(x',y_3,y_4).
\end{equation}
Substituting \cref{eq:fc3_separable_general} for both copies of
$\Phi$, factoring the $g$-independent parts, and collecting yields
\begin{align}
  \label{eq:bubble_factorised}
  \mathcal{B}_{xx'}
  &= \sum_{r,s=1}^{R} \lambda_r\,\lambda_s\;
      \alpha_r(x)\,\alpha_s(x')\;
      P_{rs}\,Q_{rs},
  \\
  \label{eq:gram_P}
  P_{rs} &\equiv \beta_r^{\!\top} g\,\gamma_s \;=\; (B^{\!\top} g C)_{rs},
  \\
  \label{eq:gram_Q}
  Q_{rs} &\equiv \gamma_r^{\!\top} g\,\beta_s \;=\; (C^{\!\top} g B)_{rs}.
\end{align}
The four-cycle in the original ring collapses to two independent
two-cycles. Each two-cycle is a scalar bilinear form on the
propagator. The full $\NBS\times\NBS$ propagator entered
\cref{eq:inner_contraction} four times through tensor contractions. In
\cref{eq:bubble_factorised,eq:gram_P,eq:gram_Q}, it enters twice as the
kernel of a quadratic form.

The matrices $P$ and $Q$ defined in \cref{eq:gram_P,eq:gram_Q} are
not independent. They satisfy the algebraic identity
% REVIEW: reference identity -- not used downstream (orphan); kept for completeness
\begin{equation}\label{eq:Q_transpose}
  Q \;=\; C^{\!\top} g\,B \;=\; (B^{\!\top} g^{\!\top} C)^{\!\top}.
\end{equation}

\emph{INDSCAL ($B=C$).} The two bilinear forms are
\begin{equation}\label{eq:P_indscal}
  P_{rs} \;=\; \beta_r^{\!\top} g\,\beta_s
         \;=\; (B^{\!\top} g B)_{rs}
         \;=\; Q_{rs},
\end{equation}
identical entry-by-entry. The result holds without any symmetry
assumption on $g$.

\emph{Waring ($A=B=C=U$).} Same as INDSCAL, with the additional
property that the assembly factor $\alpha$ also reduces to $U$:
\begin{equation}\label{eq:bubble_waring}
  \mathcal{B}_{xx'}^{\mathrm{Waring}}
  = \sum_{r,s} \lambda_r \lambda_s\,U_{x,r}\,U_{x',s}\,
      \bigl[(U^{\!\top} g\,U)_{rs}\bigr]^{2}.
\end{equation}
The bilinear form is squared elementwise; the entire kernel is
parametrised by a single $\ndof\times R$ matrix $U$ and the weights
$\lambda$.

For CP without symmetry, $P\ne Q$ in general.

\paragraph{Unified kernel.}
Under the simplifications above, the bubble \cref{eq:bubble_factorised}
can be written as a single GEMM--Hadamard--GEMM sequence:
\begin{align}
  \label{eq:kernel_step1}
  V       &= g\,W,                          \\
  \label{eq:kernel_step2}
  P       &= W^{\!\top} V,                     \\
  \label{eq:kernel_step3}
  H_{rs}  &= \lambda_r\,\lambda_s\,P_{rs}\,P^{\star}_{rs}, \\
  \label{eq:kernel_step4}
  \Sigma  &= A\,H\,A^{\!\top},
\end{align}
where the choice of $W$ and $P^{\star}$ depends on the ansatz:
\begin{equation}\label{eq:kernel_choices}
\begin{aligned}
  \text{Waring:}\;\; & W = U,\;\; P^\star = P, \;\; A = U, \\
  \text{INDSCAL:}\;\; & W = B,\;\; P^\star = P, \;\; A = A_{\mathrm{INDSCAL}}, \\
  \text{CP:}\;\; & W = (B,C),\;\; P^\star = Q, \;\; A = A_{\mathrm{CP}}.
\end{aligned}
\end{equation}
For CP, $W=(B,C)$ denotes that two SpMMs are needed (one against $B$,
one against $C$) to obtain $V_B=gB$ and $V_C=gC$. INDSCAL and Waring need only
one SpMM. This is the dominant cost saving from imposing internal
symmetry on the decomposition.

% ----------------------------------------------------------------------
\subsection{Specialisation to each ansatz}
\label{ssec:phph_decomp_ansatze}
% ----------------------------------------------------------------------

The CP-family ans\"atze all fit the unified kernel
\cref{eq:kernel_step1,eq:kernel_step2,eq:kernel_step3,eq:kernel_step4}
with the choices \cref{eq:kernel_choices}. Tucker and the
matricisation SVD require minor modifications, given below.

\paragraph{CP.}
$\Phi_{xy_1y_2}^{\mathrm{CP}} = \sum_{r} \lambda_r\,a_r(x)\,b_r(y_1)\,c_r(y_2)$
with $A,B,C$ independent. The kernel requires both $V_B = gB$ and
$V_C = gC$; $P = B^{\!\top}V_C$ and (assuming $g=g^{\!\top}$) $Q=P^{\!\top}$,
giving $H_{rs}=\lambda_r\lambda_s P_{rs}P_{sr}$. The reconstructed
tensor is generically not $S_2$-symmetric in $(y_1,y_2)$; this
asymmetry survives into $\Sigma$ unless restored in post-processing.

\paragraph{INDSCAL.}
$\Phi_{xy_1y_2}^{\mathrm{INDSCAL}} = \sum_{r}\,d_r(x)\,v_r(y_1)\,v_r(y_2)$
with $B=C\equiv V$ and $A_\mathrm{INDSCAL} = D = [d_1\,\dots\,d_R]$.
The kernel computes $V_B = gV$, $P=V^{\!\top}V_B$ (a single GEMM-like
product), and assembles via $\Sigma = D\,(P\circ P)\,D^{\!\top}$ with
the elementwise square $(P\circ P)_{rs}=\lambda_r\lambda_s P_{rs}^{2}$.
% REVIEW: S_2-symmetric reconstruction -> Sigma_{xx'}=Sigma_{x'x} requires g=g^T (equilibrium/reciprocal G)
The reconstructed tensor is exactly $S_2$-symmetric, so the SSE
inherits exact $\Sigma_{xx'}=\Sigma_{x'x}$ (requires $g=g^{\!\top}$,
i.e.\ equilibrium / reciprocal $G$) at the algorithmic level.

\paragraph{Waring.}
$\Phi_{xy_1y_2}^{\mathrm{Waring}} = \sum_{r}\lambda_r\,u_r(p(x))\,u_r(y_1)\,u_r(y_2)$,
where $p$ is the supercell lift of \cref{sub:fc3_symcp}. With
$A=B=C=U$, the kernel reduces further: a single matrix $U$
parametrises both the SpMM input and the output assembly. The bubble
% REVIEW: P=P^T (hence the SYRK) requires g=g^T (equilibrium/reciprocal G)
is symmetric in $(r,s)$ by construction, so $P$ is symmetric
($P=P^{\!\top}$, requires $g=g^{\!\top}$, i.e.\ equilibrium / reciprocal $G$)
and may be computed as a SYRK (upper-triangular GEMM)
at half the cost of the general INDSCAL GEMM.

\paragraph{PCP.}
$\Phi^{\mathrm{PCP}}_{xy_1y_2} = \sum_{\xi}\tfrac{\lambda_\xi}{6}
\sum_{\sigma\in S_3}
u^{\xi}_{\sigma_1}(p(x))\,u^{\xi}_{\sigma_2}(y_1)\,u^{\xi}_{\sigma_3}(y_2)$.
Each $\sigma\in S_3$ contributes a CP-shaped term, so PCP is
algebraically equivalent to CP at rank $6R$ with a specific structure
on the $6R$ triplets. The kernel
\cref{eq:kernel_step1,eq:kernel_step2,eq:kernel_step3,eq:kernel_step4}
applies with $R \to 6R$. For SSE evaluation, PCP is dominated by
Waring at matched fitting cost; we include it only for completeness.

\paragraph{Tucker.}
$\Phi_{xy_1y_2}^{\mathrm{Tucker}}=\sum_{pqr}G_{pqr}A_{xp}B_{y_1 q}B_{y_2 r}$
with $A\in\mathbb{R}^{\ndof\times R_1}$, $B\in\mathbb{R}^{\ndof\times R_2}$
columnwise orthonormal and a core $G\in\mathbb{R}^{R_1\times R_2\times R_2}$
with $G_{pqr}=G_{prq}$. The factorised kernel becomes matrix-valued:
\begin{align}
  \label{eq:tucker_kernel_step1}
  V_B    &= g\,B \in\mathbb{R}^{\ndof\times R_2},\\
  \label{eq:tucker_kernel_step2}
  M      &= B^{\!\top} V_B \in\mathbb{R}^{R_2\times R_2},\\
  \label{eq:tucker_kernel_step3}
  H_{pp'} &= \sum_{q r q' r'} G_{p q r}\,G_{p' q' r'}\,M_{q r'}\,M_{r q'},\\
  \label{eq:tucker_kernel_step4}
  \Sigma &= A\,H\,A^{\!\top},
\end{align}
The bilinear form is a single $R_2\times R_2$ Gram matrix $M$ (a
single SpMM-GEMM); the core contraction
\cref{eq:tucker_kernel_step3} is purely propagator-independent and
has cost $\cO(R_1^{2}R_2^{3})$ after $S_2$ symmetrisation. The
% REVIEW: cost mixes per-block (NBS) and per-tau (ndof=NB*NBS) conventions vs the CP family
propagator-side cost is $\cO(R_2\,\NBS\,r_c\,C_a + R_2^{2}\NBS)$,
identical in form to the CP family with $R\to R_2$.

% ----------------------------------------------------------------------
\subsection{Block sparsity}
\label{ssec:phph_decomp_blocks}
% ----------------------------------------------------------------------

The unified kernel
\cref{eq:kernel_step1,eq:kernel_step2,eq:kernel_step3,eq:kernel_step4}
contracts $g$ against a dense factor matrix $W$ in step
\cref{eq:kernel_step1} and contracts $H$ against $A$ in step
\cref{eq:kernel_step4}. Two of the four block-sparsity layers of
\cref{ssec:phph_sparsity} survive the decomposition unchanged; one is
absorbed into the factor vectors; one becomes the support pattern of
$P$.

\paragraph{(a) BT-band of $g$ — preserved.}
The cut-off $r_c$ on the dynamical matrix gives $g(\tau)$
a block-banded structure of bandwidth $r_c$ along the transport
direction (\cref{eq:G_bb_nnz}). The SpMM $V=g\,W$ is therefore a
block-banded matrix times a dense $\ndof\times R$ matrix. The
non-zero pattern of $V$ is dictated by $g$: for any block $I$ on the
transport axis, $V_I = \sum_{|J-I|\le r_c}\,g_{IJ}\,W_J$, a band of
width $2r_c+1$ in the input. The non-zero count of $V$ is
\begin{equation}\label{eq:V_nnz}
  \mathrm{nnz}(V) \;=\; R\cdot\mathrm{nnz}_{\text{cols}}(g)
                 \;=\; \cO\!\bigl(R\,\NBS\,\NB\bigr)
                 \;=\; \cO\!\bigl(R\,\ndof\bigr),
\end{equation}
provided $W$ is dense (factor vectors with full support along
transport). The cost of the SpMM is
$\cO(R\cdot\mathrm{nnz}(g))=\cO(R\,\ndof\,r_c\,C_a)$.

\paragraph{(b) Intra-block atom-pair sparsity of $g$ — preserved.}
Within each $\NBS\times\NBS$ block, $g_{IJ}$ has nominal density
$C_a/\NBS$ (with $C_a=\NBS$ in the present dense-block kernel; see
the forward note in \cref{ssec:phph_sparsity}). The block-level SpMM
$V_I=\sum_J g_{IJ}\,W_J$ is therefore a sparse--dense block product
with cost
$\cO(R\,\NBS\,C_a)$ per source block $J$, summed over the $\cO(r_c)$
blocks within bandwidth, giving $\cO(R\,\NBS\,r_c\,C_a)$ per output
block and $\cO(R\,\ndof\,r_c\,C_a)$ for the full SpMM.

\paragraph{(c) FC3 inter-block support — absorbed.}
The FC3 cut-off
\cref{eq:fc3_block_support} restricts $\Phi^{(3)}_{I,K_1,K_2}$ to
block triples with mutual transport-direction distance
$\le N_\mathrm{c,x}$. After decomposition, the FC3 cut-off is
absorbed into the support of the factor vectors. Specifically, for
a CP fit of the full FC3 tensor, each rank-one component
$\lambda_r\,a_r\otimes b_r\otimes c_r$ has full support
$\mathrm{supp}(a_r)\times\mathrm{supp}(b_r)\times\mathrm{supp}(c_r)$;
all such supports must lie within the FC3 cut-off, so each individual
factor vector has support within $\cO(N_\mathrm{c,x}+\mathrm{spread})$
transport blocks. The total non-zero count of $W$ is therefore
\begin{equation}\label{eq:W_nnz}
  \mathrm{nnz}(W) \;=\; \cO\!\bigl(R\,\NBS\,N_\mathrm{c,x}\bigr)
                 \;=\; \cO\!\bigl(R\,\ndof\,N_\mathrm{c,x}/\NB\bigr),
\end{equation}
substantially smaller than the dense $R\,\ndof$ when
$N_\mathrm{c,x}\ll\NB$ (true for any transport device longer than
the FC3 cut-off).

In practice, an unconstrained CP-ALS fit produces factor vectors
without compact support: the rotational gauge of CP allows globally
delocalised factors to represent locally supported tensors. To
preserve the FC3 cut-off as factor-vector support, one of three
strategies is used: (i) initialise CP-ALS from a spatially localised
seed (e.g.\ the slice-wise SVD initialisation of
\cref{sub:fc3_indscal}); (ii) regularise with an $\ell_1$ or
group-sparse penalty as in
\cite{papalexakis_parcube_2012, kim_sparse_2007}; or (iii) work in
the CANDELINC framework of \cref{sub:fc3_other} with a basis that
respects the FC3 cut-off. All three preserve \cref{eq:W_nnz} at
matched accuracy; the choice is empirical and beyond the scope of
the SSE kernel.

\paragraph{(d) Output BT band — preserved by output projection.}
The output $\Sigma$ lives on the same BT band as $G$, with
$|I-J|\le 1$. The assembly step \cref{eq:kernel_step4} is therefore
restricted to this band. \emph{Crucially}, this restriction is
applied at the output, not at the intermediate $H$: the $R\times R$
matrix $H$ has dense support and does not inherit the band structure
of $\Sigma$. We compute $H$ once per $\tau$ and project to the BT
band when constructing $\Sigma$.

\paragraph{Effective sparsity of $P$.}
The Gram matrix $P$ has entries $P_{rs} = \beta_r^{\!\top} g\,\gamma_s$.
Two effects can render $P$ sparse:
\begin{itemize}
  \item \emph{Disjoint factor supports.} If
    $\mathrm{supp}(\beta_r)\cap N_{r_c}(\mathrm{supp}(\gamma_s))=\emptyset$
    (where $N_{r_c}$ is the $r_c$-block neighbourhood induced by $g$'s
    bandwidth), then $P_{rs}=0$ exactly.
  \item \emph{Approximate orthogonality.} For Waring with
    fitted $u_r$ that are approximately orthonormal in the metric
    $g$, the off-diagonal entries of $P$ are small. This is a
    refit-dependent property and we do not assume it.
\end{itemize}
The first effect can be substantial: for FC3 with $N_\mathrm{c,x}=3$
and a device of $\NB$ transport blocks, the fraction of $P_{rs}$
pairs with overlapping supports scales as
$N_\mathrm{c,x}/\NB$ — between $10^{-2}$ and $10^{-1}$ at production
device sizes. Exploiting this sparsity in step
\cref{eq:kernel_step2} converts the dense $R^{2}\NBS$ GEMM into a
sparse--sparse product at cost $\cO(R^{2}N_\mathrm{c,x})$, dominated
by the SpMM cost.

% ----------------------------------------------------------------------
\subsection{Computation strategy}
\label{ssec:phph_decomp_compute}
% ----------------------------------------------------------------------

The per-$\tau$ kernel
\cref{eq:kernel_step1,eq:kernel_step2,eq:kernel_step3,eq:kernel_step4}
maps onto three numerical primitives, all standard in modern linear
algebra libraries:
\begin{enumerate}
  \item \emph{Step \cref{eq:kernel_step1}: batched sparse SpMM.}
        $V = g\,W$, where $g$ is block-banded sparse and $W$ is
        $\ndof\times R$ with FC3-induced support \cref{eq:W_nnz}.
        Implemented as a sparse matrix times dense matrix product
        with $R$ right-hand sides batched into a single SpMM call.
        Cost: $\cO(R\,\ndof\,r_c\,C_a)$ flops, memory traffic
        $\cO(\mathrm{nnz}(g) + R\,\ndof)$. With $R$ batched RHSs,
        the arithmetic intensity is $\cO(R)$ flops/byte, placing the
        kernel firmly in the compute-bound regime of modern
        accelerators for $R\gtrsim 32$.
  \item \emph{Step \cref{eq:kernel_step2}: small dense GEMM.}
        $P = W^{\!\top}V$, a dense $R\times\ndof$ times $\ndof\times R$
        product yielding an $R\times R$ matrix.
        For Waring, $P$ is symmetric and the product reduces to
        SYRK. Cost: $\cO(R^{2}\,\ndof)$ flops. With $R^{2}$ output
        entries and $\cO(R\,\ndof)$ input non-zeros, the arithmetic
        % REVIEW: O(ndof/R) intensity is inconsistent with the standard GEMM roofline O(R) and the compute-bound claim
        intensity is $\cO(\ndof/R)$ flops/byte, comfortably
        compute-bound.
  \item \emph{Step \cref{eq:kernel_step4}: sparse-output sandwich.}
        $\Sigma = A\,H\,A^{\!\top}$ restricted to the output BT band
        and any available intra-block atom-pair sparsity
        (with $C_a=\NBS$ in the present dense-block kernel; see the
        forward note in \cref{ssec:phph_sparsity}).
        Implemented as: (i) precompute $T=AH$, a dense $\ndof\times R$
        matrix, cost $\cO(R^{2}\ndof)$; (ii) for each non-zero
        entry $(x,x')$ in the output support, compute
        $\Sigma_{xx'}=T_{x,:}\cdot A_{x',:}^{\!\top}$, cost
        $\cO(R)$. Total: $\cO(R^{2}\ndof + R\,\ndof\,C_a)$.
\end{enumerate}
The Hadamard step \cref{eq:kernel_step3} has cost $\cO(R^{2})$ and
is negligible.

\paragraph{Per-$\tau$ cost.}
Summing the dominant terms,
\begin{equation}\label{eq:per_tau_cost}
  \mathcal{C}_{\tau}^{\mathrm{decomp}}
  \;=\;
  \underbrace{\cO\!\bigl(R\,\ndof\,r_c\,C_a\bigr)}_{\text{SpMM}}
  \;+\;
  \underbrace{\cO\!\bigl(R^{2}\,\ndof\bigr)}_{\text{Gram + assembly}}.
\end{equation}
This is to be compared with the sparse-ring per-$\tau$ cost
\cref{eq:cost_phph_total},
\begin{equation}\label{eq:per_tau_cost_sparse}
  \mathcal{C}_{\tau}^{\mathrm{sparse}}
  \;=\;
  \cO\!\bigl(\ndof\,N_\mathrm{c,x}^{2}\,C_a^{3}\bigr).
\end{equation}
At Si-nanoribbon parameters ($C_a\sim 30$, $N_\mathrm{c,x}=3$,
$r_c\sim 2$), the sparse-ring cost has prefactor
$N_\mathrm{c,x}^{2}C_a^{3}\sim 2.4\times 10^{5}$; the factorised
kernel has prefactor $R\,r_c\,C_a + R^{2}\sim 60R + R^{2}$. For
% REVIEW: 2.4e5/1.6e4 ~ 15x is about one order of magnitude, not two
$R=100$, the factorised kernel is roughly an order of magnitude
($\sim 15\times$) cheaper.

\paragraph{Crossover.}
The factorised kernel is cheaper than the sparse ring when
\begin{equation}\label{eq:crossover}
  R\,r_c\,C_a + R^{2}
  \;<\; N_\mathrm{c,x}^{2}\,C_a^{3},
\end{equation}
which, for $R\ll C_a$, simplifies to $R\lesssim N_\mathrm{c,x}^{2}\,C_a^{2}/r_c\sim 10^{4}$
at the same parameters; for $R$ comparable to $C_a$, the $R^{2}$
term enters and the crossover moves to
$R\lesssim N_\mathrm{c,x}\,C_a^{3/2}\sim 500$. Reported CP and
Waring ranks for FC3 of Si, MgO, HgTe, TiNiSn, ZrO$_2$
\cite{luo_tensor_2025} are well below these bounds, placing the
factorised kernel firmly in the favourable regime.

\paragraph{Batching over $\tau$.}
The kernel is local in $\tau$: each $\tau$-slice produces an
independent $\Sigma(\tau)$. The natural batching strategy is to
process $N_b$ $\tau$-slices in a single call, stacking the
$N_b$ propagator slices into a $\ndof\times\ndof\times N_b$ tensor
and the corresponding $V$ matrices into a $\ndof\times R\times N_b$
tensor. The SpMM becomes a batched SpMM (cuSPARSE/MKL support);
the GEMM becomes a batched GEMM. Batch size $N_b$ is chosen
adaptively from the free HBM, with the dominant memory term being
the $V$ tensor at $\cO(N_b R\,\ndof)$.

\paragraph{Sparse-factor optimisation.}
When the factor vectors $W$ are intra-block sparse with density
$C_a/\NBS$ (the case for FC3-localised factors, see
\cref{ssec:phph_decomp_blocks}), the GEMM in
\cref{eq:kernel_step2} can be replaced by a sparse--dense product at
cost $\cO(R^{2}\,C_a\,N_\mathrm{c,x})$, reducing the second term in
\cref{eq:per_tau_cost} by a factor $\NBS/(C_a N_\mathrm{c,x})$. The
SpMM cost is unaffected (it is already proportional to
$\mathrm{nnz}(g)$, not to the density of $W$).

% ----------------------------------------------------------------------
\subsection{Transverse periodic case}
\label{ssec:phph_decomp_qspace}
% ----------------------------------------------------------------------

In the transverse-periodic setting of \cref{sub:transper_negf}, the
SSE \cref{eq:sigma_phph} carries an internal sum over a transverse
wavevector $\qp'$ which couples $g$ at $\qp'$ and $\qp-\qp'$ via the
FC3 evaluated at the same pair of momenta. We carry the
factorised kernel through the partial Fourier transform and identify
the structural simplifications.

\paragraph{Partial Fourier transform of CP factors.}
Performing the CP decomposition on the real-space FC3
$\Phi^{(3)}_{\mu,\,l_1\nu_1,\,l_2\nu_2}$ gives factor matrices in
which the second and third legs include the transverse cell label.
The partial Fourier transform \cref{eq:fc3_fourier} acts on these
cell labels and produces
\begin{equation}\label{eq:cp_fourier_decomp}
  \tilde\Phi^{(3)}_{\mu\nu_1\nu_2}(\qp,\qp')
  \;=\; \sum_{r}\lambda_r\,a_r(\mu)\,b_r(\nu_1;\qp)\,c_r(\nu_2;\qp'),
\end{equation}
with $\qp$-dependent factor vectors
$b_r(\nu_1;\qp)=\sum_{l_1}\tilde b_r(l_1,\nu_1)e^{-i\qp\cdot l_1}$ and
likewise for $c_r$. The CP rank is preserved exactly under the
partial Fourier transform: the decomposition
\cref{eq:fc3_separable_general} commutes with the cell-FT because the
FT acts on each factor leg separately. For Waring, the single
factor vector $u_r$ has all three legs Fourier-transformed against
the appropriate cell labels.

\paragraph{Factorised kernel in $\qp$-space.}
With $\Phi_L=\Phi_R$ at $(\qp',\qp-\qp')$ and $(\qp'-\qp,-\qp')$
respectively, and noting that for real factor vectors
$b_r(\nu;-\qp)=b_r(\nu;\qp)^{*}$, the kernel
\cref{eq:bubble_factorised} becomes
\begin{align}
  \label{eq:bubble_qspace}
  \mathcal{B}_{xx'}(\tau,\qp,\qp')
  &= \sum_{rs}\lambda_r\lambda_s\,a_r(x)\,a_s(x')\,
      P_{rs}(\tau,\qp,\qp')\,Q_{rs}(\tau,\qp,\qp'),\\
  % REVIEW: conjugated factor sits on the g(q') line -> momentum q' (not q-q'), conserving momentum
  \label{eq:P_qspace}
  P_{rs}(\tau,\qp,\qp')
  &= b_r(\cdot;\qp')^{\!\top}\,g(\tau,\qp')\,c_s(\cdot;\qp')^{*},\\
  % REVIEW: conjugated factor sits on the g(q-q') line -> momentum q-q' (= q'-q*), conserving momentum
  \label{eq:Q_qspace}
  Q_{rs}(\tau,\qp,\qp')
  &= c_r(\cdot;\qp-\qp')^{\!\top}\,g(\tau,\qp-\qp')\,b_s(\cdot;\qp-\qp')^{*}.
\end{align}
The two bilinear forms now involve $g$ at \emph{different} transverse
momenta, $\qp'$ and $\qp-\qp'$, and are no longer related by transposition
even for Waring. They are, however, related by a change of summation
variable: under $\qp'\to\qp-\qp'$,
$P_{rs}(\tau,\qp,\qp-\qp')=Q_{rs}(\tau,\qp,\qp')$. After summing over
the full Monkhorst--Pack mesh (which is closed under
$\qp'\to\qp-\qp'$ for appropriate shifts), the contributions of $P$ and
$Q$ to $\Sigma$ are equal, so only $P$ need be evaluated.

\paragraph{Kernel restructure.}
The change-of-variable identity collapses
\cref{eq:bubble_qspace,eq:P_qspace,eq:Q_qspace} to
\begin{equation}\label{eq:sigma_qspace_final}
  % REVIEW: restore i\hbar/2 prefactor of eq:sigma_phph; drop spurious factor 2 (full q'-sum already double-counts P<->Q)
  \Sigma_{xx'}(\tau,\qp)
  = \frac{i\hbar}{2N}\!\sum_{\qp'}\!
    \sum_{rs}\lambda_r\lambda_s a_r(x)a_s(x')\,
    P_{rs}(\tau,\qp,\qp')\,P_{rs}(\tau,\qp,\qp-\qp'),
\end{equation}
which is a convolution along $\qp'$ of the matrix-valued quantity
$P(\tau,\qp,\qp')$ with its shift $P(\tau,\qp,\qp-\qp')$.
% REVIEW: "single FFT" understates it -- the q'-convolution needs FFTs over the nu/r/s structure
The
$\qp'$-sum is therefore a Fourier transform of a per-block product, and
the $N_q^{2}$ factor in \cref{eq:cost_q} is replaced by
$N_q \log N_q$ after a single FFT along the transverse axis.
% REVIEW: full q'-sum already includes both orderings, so no extra factor 2 (would apply only under half-BZ folding)
Summing over the full $\qp'$ mesh already accounts for both orderings
$P(\qp')P(\qp-\qp')$, so no extra factor of 2 appears; for CP without
internal symmetry $P$ and $Q$ differ and both bilinear forms enter
explicitly.

% ----------------------------------------------------------------------
\subsection{Distributed implementation}
\label{ssec:phph_decomp_dist}
% ----------------------------------------------------------------------

The four-redistribution pipeline of \cref{ssec:phph_distribution} is
preserved bit-for-bit. The decomposition changes only Stage~2 of the
pipeline (the per-$\tau$ ring contraction). Stages~1 and~3
(FFT $G\to g$ and IFFT $\sigma\to\Sigma$), the two interconvertible
distributions \texttt{stack}/\texttt{nnz}, and the four
\texttt{stack}$\leftrightarrow$\texttt{nnz} dtransposes are unchanged.

\paragraph{Stage~2 with decomposition.}
The ring contraction
\cref{eq:phph_step_A,eq:phph_step_B,eq:phph_step_S} is replaced by
the unified kernel
\cref{eq:kernel_step1,eq:kernel_step2,eq:kernel_step3,eq:kernel_step4}.
On a single rank holding a $\tau$-slice of $g$ in the \texttt{stack}
distribution, the kernel is purely local in $\tau$ and consists of
SpMM, GEMM, and the sparse-output sandwich. No communication is
incurred inside Stage~2 except the FC3-dictionary halo discussed
below.

\paragraph{Per-rank arithmetic.}
Summing over $N_\tau$, $N_q^{2}$ (or $N_q\log N_q$ in the
$\qp'\to\qp-\qp'$ symmetrised form of
\cref{eq:sigma_qspace_final}) and distributing across $P=P_E\,P_S$
ranks gives
\begin{equation}\label{eq:cost_phph_dist_cp}
  W^{\mathrm{rank}}_\mathrm{phph,decomp}
  \;=\;
  \cO\!\Bigl(
    \tfrac{N_\tau\,N_q^{2}\,N_{AO}\,(R\,r_c\,C_a + R^{2})}{P}
  \Bigr).
\end{equation}
% REVIEW: 2.4e5 = N_cx^2 C_a^3 is the sparse baseline (cost_phph_total), not the dense-block cost_phph_dist
The $C_a^{3}\,N_\mathrm{c,x}^{2}$ factor of \cref{eq:cost_phph_total}
is replaced by $R\,r_c\,C_a + R^{2}$. At $C_a=30$, $r_c=2$, $R=100$:
the old factor is $2.4\times 10^{5}$, the new factor is
$\sim 1.6\times 10^{4}$ --- a 15$\times$ arithmetic reduction. With
$\qp'$-FFT symmetrisation (\cref{ssec:phph_decomp_qspace}), an
% REVIEW: N_q=400 here vs N_q=64 (8x8) in the streaming example below -- two different example meshes
additional factor $N_q/\log N_q\sim 50$ is gained for production
$N_q=400$.
 
% ----------------------------------------------------------------------
\subsection{Transverse-momentum (q-point) distribution and vertex streaming}
\label{ssec:phph_qdist}
% ----------------------------------------------------------------------

The distribution of \cref{ssec:phph_distribution,ssec:phph_decomp_dist}
spreads the work over frequency ($P_E$, \texttt{comm.stack}) and
transport blocks ($P_S$, \texttt{comm.block}). In the
transverse-periodic setting of \cref{ssec:phph_decomp_qspace} a third
index appears: the external transverse momentum $\qp$. Unlike the
finite (zone-centre) device, where each $\qp$ is an independent
external label, the periodic self-energy
\begin{equation}\label{eq:sigma_q_dist}
  \Sigma^{\lessgtr}(\qp,\omega)
  = \frac{i\hbar}{2}\frac{1}{N_q}\sum_{\qp'}\!\int\!\frac{d\omega'}{2\pi}\,
    \tilde\Phi(\qp',\qp{-}\qp')\,
    G^{\lessgtr}(\qp',\omega')\,G^{\lessgtr}(\qp{-}\qp',\omega{-}\omega')\,
    % REVIEW: use established convention Phi~(q'-q,-q') of eq:sigma_phph (differs by a reality conjugation)
    \tilde\Phi(\qp'{-}\qp,-\qp')
\end{equation}
\emph{couples} the momenta through crystal-momentum conservation
$\qp=\qp'+(\qp-\qp')$: every external $\qp$ draws on the Green's
functions of \emph{all} internal $\qp'$. The momenta are therefore not
embarrassingly parallel, and a naive per-$\qp$ split would replicate
the full $G(\qp')$ set on every rank.

\paragraph{The q-communicator.}
We add $\qp$ as a third communicator axis alongside block and stack, so
that the global ranks factor as
$q\times\mathrm{stack}\times\mathrm{block}$. Concretely
\texttt{QuatrexCommunicator.configure} gains a \texttt{q\_comm\_size}
and a third \texttt{MPI\_Comm} split; the previous two-axis layout is
recovered exactly at \texttt{q\_comm\_size}~$=1$ (the $q$-communicator
then has one rank and is a no-op), so block and stack behaviour is
unchanged. Each $q$-rank owns a disjoint slice of the \emph{external}
$\qp$ mesh and stores only its own per-$\qp$ Green's functions and
output self-energies. The \emph{internal} $\qp'$ Green's functions that
\cref{eq:sigma_q_dist} needs are reconstructed once per SCBA iteration
by an \texttt{all\_gather\_v} over the $q$-communicator -- the
``data exchange between processors'' inherent to the momentum
convolution. The explicit external-$\qp$ loop then divides as $1/P_q$,
while the all-gather of the per-$\qp$ diagonal $G$ carries the
communication cost; the bubble kernel itself is unchanged. This third
axis composes with the energy and block axes without touching them: a
production run can use all three simultaneously
($P=P_q\,P_E\,P_S$).

\paragraph{Streaming the momentum-space vertex.}
On an accelerator the binding constraint for
\cref{eq:sigma_q_dist} is memory rather than arithmetic. The
Fourier-transformed FC3 vertex
$\tilde\Phi(\qp,\qp')=T(\qp)\,M\,T(\qp')^{\dagger}$ -- with $T(\qp)$
the supercell-to-primitive gathering matrix and $M$ the (small,
real-space) mass-weighted FC3 -- is a dense rank-five object of size
$N_q^{2}\,\ndof^{3}$ if materialised. At an $8\times8$ transverse mesh
($N_q=64$) and a wire-scale cross section ($\ndof=63$) this is already
$\sim 16$~GB, beyond a single device, and it grows as $N_q^{2}$. Since
the kernel only ever touches the pair
$(\tilde\Phi(\qp',\qp{-}\qp'),\,\tilde\Phi(\qp{-}\qp',\qp'))$ for the
momenta in the current batch, we never form the full tensor: with
\texttt{stream\_phi} enabled the kernel keeps only $T$ and $M$ and
rebuilds $\tilde\Phi$ on the fly per momentum batch, dropping the peak
from $\cO(N_q^{2}\ndof^{3})$ to $\cO(N_q\,\ndof^{2})$. This trades a
vertex recomputation for memory, which is favourable on a GPU where the
bandwidth-bound bubble contractions dominate the run time regardless,
and it composes with the $q$-communicator at no extra communication
(each rank streams only the vertices for its own external-$\qp$ slice).
The result is algebraically identical to the dense path; the measured
memory reduction and the near-ideal $q$-axis speed-up are reported in
\cref{sec:results}.
